\section{Applicable standards and how they are identified}
To identify the applicable harmonised EMC emission standards, the European Commission
webpage for harmonised standards under the EMC Directive 2014/30/EU was used as the
starting point. This page provides a ``Summary list of titles and references of harmonised
standards under Directive 2014/30/EU for EMC'', available as a downloadable PDF or
spreadsheet. The summary list contains, among other information, the reference number of
each standard and the title of the standard.

The list was first reviewed by comparing the standard titles with the intended environment
and function of the AGV.
some candidates are:
The following candidate standards were selected from the summary list for further
comparison:

\begin{table}[H]
\centering
\caption{Candidate harmonised EMC standards selected from the summary list under Directive 2014/30/EU.}
\small
\setlength{\tabcolsep}{4pt}
\begin{tabular}{@{}p{0.30\linewidth}p{0.32\linewidth}p{0.32\linewidth}@{}}
\hline
{\raggedright\textbf{Standard}\par} & {\raggedright\textbf{Title in the summary list}\par} & {\raggedright\textbf{Reason for inclusion}\par} \\
\hline
{\raggedright EN~61800-3:2004+A1:2012\par} &
{\raggedright Adjustable speed electrical power drive systems, Part 3: EMC requirements and specific test methods\par} &
{\raggedright Relevant because the AGV contains motor drive electronics. It is included as a technology-related candidate, but it may only apply to the drive system rather than the complete AGV.\par} \\

{\raggedright EN~61000-6-3:2007+A1:2011\par} &
{\raggedright Electromagnetic compatibility (EMC), Part 6-3: Generic standards, Emission standard for residential, commercial and light-industrial environments\par} &
{\raggedright Relevant as a generic emission candidate because the intended environment has been classified as light industrial.\par} \\

{\raggedright EN~61000-6-4:2007+A1:2011\par} &
{\raggedright Electromagnetic compatibility (EMC) - Part 6-4: Generic standards - Emission standard for industrial environments\par} &
{\raggedright Relevant as the second generic emission candidate. It is included for comparison because the AGV operates in a technical automated environment, although the environment is probably not treated as heavy industrial.\par} \\

{\raggedright EN~14010:2003+A1:2009\par} &
{\raggedright Safety of machinery, Equipment for power driven parking of motor vehicles, Safety and EMC requirements for design, manufacturing, erection and commissioning stages\par} &
{\raggedright Relevant as the strongest product-related candidate because the AGV is intended for automated parking of motor vehicles.\par} \\
\hline
\end{tabular}
\end{table}

The best fit is EN~14010:2003+A1:2009, \textit{Safety of machinery --- Equipment for power driven parking of motor vehicles --- Safety and EMC requirements for design, manufacturing, erection and commissioning stages}.

As this standard is lincence based it cant be refrence direclty but EN 14010 does not itself conatain EMC standard instead it stats that: "The electromagnetic disturbances generated by the power driven parking equipment shall not exceed the levels
specified in generic emission standard EN 61000-6-3"

There for focus will be on EN 61000-6-3, \textit{Electromagnetic compatibility (EMC), Part 6-3: Generic standards, Emission standard for residential, commercial and light-industrial environments}.

Since EN 61000-6-3 standard also is licence-based, this test report does not claim
formal conformity with the standard. The standard cant be used directly, instead it is used only as a practical
pre-compliance reference, and the numerical limits used in the measurements are taken
from publicly available secondary sources that reference EN 61000-6-3.


\begin{table}[H]
\centering
\caption{Secondary sources used for practical radiated-emission reference limits.}
\small
\setlength{\tabcolsep}{4pt}
\begin{tabular}{@{}p{0.25\linewidth}p{0.24\linewidth}p{0.16\linewidth}p{0.27\linewidth}@{}}
\hline
{\raggedright\textbf{Frequency range}\par} &
{\raggedright\textbf{Limit}\par} &
{\raggedright\textbf{Page}\par} &
{\raggedright\textbf{Source}\par} \\
\hline

{\raggedright \(30\,MHz\) to \(230\,MHz\)\par} &
{\raggedright \(40\,dB\mu V/m\) QP at \(3\,m\)\par} &
{\raggedright p.~8 of 20\par} &
{\raggedright \href{https://www.steinertsensingsystems.com/wp-content/uploads/2013/06/Certificate-of-Compliance-CE-FCC-SuperChrono-1.pdf}{BK Services, EMC Test Report SuperChrono, EN~61000-6-3:2007, Radiated Emission}\par} \\

{\raggedright \(230\,MHz\) to \(1000\,MHz\)\par} &
{\raggedright \(47\,dB\mu V/m\) QP at \(3\,m\)\par} &
{\raggedright p.~8 of 20\par} &
{\raggedright \href{https://www.steinertsensingsystems.com/wp-content/uploads/2013/06/Certificate-of-Compliance-CE-FCC-SuperChrono-1.pdf}{BK Services, EMC Test Report SuperChrono, EN~61000-6-3:2007, Radiated Emission}\par} \\

{\raggedright \(30\,MHz\) to \(230\,MHz\)\par} &
{\raggedright \(40\,dB\mu V/m\) QP at \(3\,m\)\par} &
{\raggedright p.~11 of 19\par} &
{\raggedright \href{https://www.sapv.fr/web/bundles/public/fichiers/7a086ddf41-caracteristiques.pdf}{CETIM, Essais de compatibilit\'e \'electromagn\'etique, NF~EN~61000-6-3, radiated emission limits}\par} \\

{\raggedright \(230\,MHz\) to \(1000\,MHz\)\par} &
{\raggedright \(47\,dB\mu V/m\) QP at \(3\,m\)\par} &
{\raggedright p.~11 of 19\par} &
{\raggedright \href{https://www.sapv.fr/web/bundles/public/fichiers/7a086ddf41-caracteristiques.pdf}{CETIM, Essais de compatibilit\'e \'electromagn\'etique, NF~EN~61000-6-3, radiated emission limits}\par} \\

\hline
\end{tabular}
\end{table}

Based on the secondary sources, it can be inferred that the practical
EN~61000-6-3 radiated-emission reference limits at \(3\,m\) are
\(40\,dB\mu V/m\) from \(30\,MHz\) to \(230\,MHz\), and
\(47\,dB\mu V/m\) from \(230\,MHz\) to \(1000\,MHz\), using a quasi-peak
detector.



